\begin{frame}{다음 주차 예고: Week 6. 시간차 학습 (TD)}
  \begin{itemize}
    \item 이번 장의 MC 는 \textbf{모델 없이} 배운다는 대가로
      \textbf{에피소드를 끝까지 기다려야} 했다.
    \item TD 는 그 기다림을 없앤다 --- \textbf{매 스텝}에서
      ``한 스텝만 보고도 즉시 갱신''
      ($Q$ 를 한 스텝 분만 ``내다보고'' 갱신).
    \item \kb{SARSA} (on-policy TD)와 \kb{Q-learning}
      (off-policy TD, 중요도 샘플링 \textbf{없이} off-policy 를
      구현)를 배우고, CliffWalking 예제로 두 방법의 차이를
      체감한다.
    \item ``모델의 유무'' × ``갱신을 언제 하는가''라는
      강화학습 지도의 두 축이 모두 놓인다 --- Week 9 DQN
      (신경망 + 경험재현)과 Week 11 PPO (클리핑된 확률비)로
      이어지는 모든 것이 이번 두 주(Week 5 MC + Week 6 TD)의
      위에서 세워진다.
  \end{itemize}
\end{frame}
